\section{Textual History of Luke 24:12}

The verse's textual history matters because removing it leaves John as the sole
canonical cloth inventory and eliminates Luke's explicit statement of Peter's
wonder. Some early witnesses omit the verse, while others include it; modern
critical judgment generally prints it with apparatus notice. The registered
Robinson--Pierpont artifact includes it and reports the modern eclectic form's
difference concerning the participle ``lying.''

Westcott and Hort argued that the verse was an assimilation from John
20:3--10, shortened by omission of the other disciple. Their judgment belongs
to a nineteenth-century textual theory and must not be reported as present
consensus. Conversely, canonical familiarity must not erase the external and
internal evidence that made the verse debatable.

Three conclusions survive either decision. First, John independently attributes
careful observation to Peter and belief to the other disciple. Second, an
included Luke 24:12 supplies early reception of Peter as wondering. Third, an
omitted verse cannot be used to make John say Peter believed. The textual
variant changes the amount of canonical evidence about Peter's immediate state;
it does not transfer John's predicate to him.

